\documentclass[11pt,a4paper]{article}
\usepackage[margin=20mm,headheight=15pt]{geometry}
\usepackage{fontspec}
\setmainfont{TeX Gyre Pagella}
\setmonofont{DejaVu Sans Mono}[Scale=0.82]
\usepackage{amsmath}
\usepackage{unicode-math}
\setmathfont{TeX Gyre Pagella Math}
\setcounter{MaxMatrixCols}{20}
\usepackage{microtype}
\usepackage{enumitem}
\usepackage{fvextra}
\usepackage{booktabs}
\usepackage{longtable}
\usepackage{fancyhdr}
\usepackage[hidelinks]{hyperref}
\setlength{\parindent}{0pt}
\setlength{\parskip}{6pt}
\setlength{\emergencystretch}{3em}
\hyphenpenalty=10000
\exhyphenpenalty=10000
\raggedright
\DefineVerbatimEnvironment{Code}{Verbatim}{fontsize=\small,
  breaklines=true,breakanywhere=true,breaksymbolleft={},breaksymbolright={},
  frame=single,framesep=5pt}
\setlist[enumerate]{leftmargin=7mm,itemsep=3pt,topsep=5pt,parsep=0pt}
\pagestyle{fancy}
\fancyhf{}
\fancyhead[L]{\small Computer Vision 1}
\fancyhead[R]{\small Exam 2 · Questions}
\fancyfoot[C]{\small \thepage}
\hypersetup{pdftitle={Exam 2 — Implementation and failure diagnosis},pdfsubject={CV1 practice material}}
\begin{document}
{\LARGE\bfseries Exam 2 — Implementation and failure diagnosis}\par
{\large Questions only}\par\bigskip
\textbf{Name:} \rule{70mm}{0.4pt}\hfill \textbf{Date:} \rule{30mm}{0.4pt}
\section*{Conditions}
\begin{itemize}[leftmargin=5mm,itemsep=5pt]
\item 50 questions: 30 single-answer multiple choice and 20 numerical. Each is worth 1 raw practice point.
\item Suggested training limit: \textbf{150 minutes}. This is not a claim about the actual exam duration or question count.
\item Use up to three double-sided A4 sheets,\allowbreak{} prepared by you,\allowbreak{} and a local Jupyter notebook. No internet,\allowbreak{} AI assistance or answer key during the attempt.
\item Use Python/\allowbreak{}NumPy for computation; SciPy is also allowed. You do not need executable OpenCV or PyTorch. Specialized routines are provided when needed.
\item For multiple choice,\allowbreak{} select exactly one of A–D. Attempt every question.
\item Unless an exact integer is requested,\allowbreak{} round the final numerical answer to \textbf{two decimal places}. Keep full precision in intermediate steps. Use fractions rather than percentages unless specified.
\item Numerical practice grading uses the syllabus rule: within ±0.01 of the unrounded reference,\allowbreak{} inclusive. Counts require the exact integer; no partial credit.
\item Conventions stated in a question override defaults. Arrays use [row,\allowbreak{}column] =\allowbreak{} [y,\allowbreak{}x]; geometric vectors are columns; image-plane projection is virtual/\allowbreak{}front-facing; foreground is 1.
\item This is original,\allowbreak{} challenging training material based on the supplied course files,\allowbreak{} not a lecturer-issued mock or a prediction of the exam.
\end{itemize}
Solutions are in a separate PDF. Do not open them until you finish.
\clearpage\section*{Part A — Single-answer multiple choice}
\begin{minipage}{\linewidth}
\textbf{Q01 — Multiple choice}\par\smallskip
\(R_1\) and \(R_2\) are proper 3D rotation matrices. Which statement is correct?
\begin{enumerate}[label=\textbf{\Alph*.}]
\item The inverse is always \(R_2^{\mathsf T}R_1^{\mathsf T}\).
\item A rotation matrix has 9 independent degrees of freedom.
\item The inverse of \(R_2 R_1\) is \(R_1^{\mathsf T}R_2^{\mathsf T}\); rotation order generally matters.
\item \(R_1 R_2\) always equals \(R_2 R_1\).
\end{enumerate}
\end{minipage}
\par\vspace{12pt}
\begin{minipage}{\linewidth}
\textbf{Q02 — Multiple choice}\par\smallskip
Crop an image by removing 40 pixels from the left and 20 from the top,\allowbreak{} with no resize. In the continuous-coordinate model,\allowbreak{} how does K change?
\begin{enumerate}[label=\textbf{\Alph*.}]
\item Subtract 40 from \(f_x\) and 20 from \(f_y\).
\item Keep K unchanged and add \(\begin{bmatrix}40&20&0\end{bmatrix}\) to camera translation.
\item Subtract 40 from \(c_x\) and 20 from \(c_y\); keep focal lengths and skew unchanged.
\item Multiply both focal lengths by the crop's remaining image area.
\end{enumerate}
\end{minipage}
\par\vspace{12pt}
\begin{minipage}{\linewidth}
\textbf{Q03 — Multiple choice}\par\smallskip
For floating-point RGB image I of shape (H,\allowbreak{}W,\allowbreak{}3),\allowbreak{} what does this code compute?


\begin{Code}
out = (I.max(axis=2) + I.min(axis=2)) / 2
\end{Code}

\begin{enumerate}[label=\textbf{\Alph*.}]
\item One scalar containing the midpoint of the global intensity range.
\item A lightness-based grayscale image of shape (H,\allowbreak{}W).
\item An RGB image with corrected white balance.
\item A hue map obtained by averaging the red and blue channels.
\end{enumerate}
\end{minipage}
\par\vspace{12pt}
\begin{minipage}{\linewidth}
\textbf{Q04 — Multiple choice}\par\smallskip
Which expression computes per-channel image means that broadcast correctly onto an array of shape (H,\allowbreak{}W,\allowbreak{}3)?
\begin{enumerate}[label=\textbf{\Alph*.}]
\item \texttt{I.mean(axis=\allowbreak{}0,\allowbreak{} keepdims=\allowbreak{}True),\allowbreak{} which always has shape (1,\allowbreak{}1,\allowbreak{}3)}
\item \texttt{I.mean(axis=\allowbreak{}(0,\allowbreak{}1),\allowbreak{} keepdims=\allowbreak{}True)}
\item \texttt{I.mean(axis=\allowbreak{}(0,\allowbreak{}1,\allowbreak{}2),\allowbreak{} keepdims=\allowbreak{}True)}
\item \texttt{I.mean(axis=\allowbreak{}2,\allowbreak{} keepdims=\allowbreak{}True)}
\end{enumerate}
\end{minipage}
\par\vspace{12pt}
\begin{minipage}{\linewidth}
\textbf{Q05 — Multiple choice}\par\smallskip
Suppose the true lighting matrix is \(L_{\mathrm{true}}=2L_{\mathrm{used}}\),\allowbreak{} with no shadows or noise. You solve \(i=L_{\mathrm{used}}\widehat g\). How do the estimated normal and albedo relate to the true ones?
\begin{enumerate}[label=\textbf{\Alph*.}]
\item The normal doubles in length after normalization; albedo is unchanged.
\item Both normal direction and albedo are unchanged.
\item The normal reverses and albedo is halved.
\item The normal is unchanged; the albedo estimate is doubled.
\end{enumerate}
\end{minipage}
\par\vspace{12pt}
\begin{minipage}{\linewidth}
\textbf{Q06 — Multiple choice}\par\smallskip
For the intrinsic-image model \(I=A\odot S\),\allowbreak{} where multiplication is elementwise,\allowbreak{} which recolouring operation preserves the estimated shading?
\begin{enumerate}[label=\textbf{\Alph*.}]
\item Divide the target albedo by S.
\item Change the object's albedo A and multiply the changed A by the original S.
\item Change S to the target RGB colour and discard A.
\item Replace every RGB pixel in I by the same target colour,\allowbreak{} including shadows.
\end{enumerate}
\end{minipage}
\par\vspace{12pt}
\begin{minipage}{\linewidth}
\textbf{Q07 — Multiple choice}\par\smallskip
A uniformly painted matte object has spatially constant albedo. Why can its image still contain a brightness gradient?
\begin{enumerate}[label=\textbf{\Alph*.}]
\item Lambertian reflection requires all pixels to have the same intensity.
\item Constant albedo forces constant depth.
\item Its surface normals and therefore light–normal dot products can vary.
\item Its intrinsic albedo must contain every shadow.
\end{enumerate}
\end{minipage}
\par\vspace{12pt}
\begin{minipage}{\linewidth}
\textbf{Q08 — Multiple choice}\par\smallskip
What information can a joint RGB histogram retain that three separate channel histograms lose?
\begin{enumerate}[label=\textbf{\Alph*.}]
\item Which red,\allowbreak{} green and blue values occur together at the same pixels.
\item The scene's illumination-independent albedo distribution,\allowbreak{} regardless of camera response.
\item The order of coloured regions across image rows,\allowbreak{} rather than just their colour counts.
\item The exact spatial location of each pixel having a particular colour.
\end{enumerate}
\end{minipage}
\par\vspace{12pt}
\begin{minipage}{\linewidth}
\textbf{Q09 — Multiple choice}\par\smallskip
A binary erosion uses a cross-shaped 3 by 3 structuring element with zeros at its corners. How should the zeros be treated?
\begin{enumerate}[label=\textbf{\Alph*.}]
\item Require the image's corner positions to be zero.
\item Replace the operation by an average over all nine pixels.
\item Include the mask zeros in a logical AND,\allowbreak{} forcing every output to zero.
\item Ignore those positions; require foreground only at active element positions.
\end{enumerate}
\end{minipage}
\par\vspace{12pt}
\begin{minipage}{\linewidth}
\textbf{Q10 — Multiple choice}\par\smallskip
A system labels every pixel by semantic class,\allowbreak{} including pixels of several touching people,\allowbreak{} but gives no person-instance IDs. What does its output provide,\allowbreak{} and can connected-component labelling alone reliably supply the missing instance separation?
\begin{enumerate}[label=\textbf{\Alph*.}]
\item Semantic segmentation; no,\allowbreak{} touching people can merge into one connected component.
\item Instance segmentation; a semantic class label automatically identifies which individual person owns each pixel.
\item Object detection only; connected-component labels are already separate bounding-box detections.
\item Panoptic segmentation; instance IDs are unnecessary as long as all pixels have semantic labels.
\end{enumerate}
\end{minipage}
\par\vspace{12pt}
\begin{minipage}{\linewidth}
\textbf{Q11 — Multiple choice}\par\smallskip
Why is a median filter often preferable to averaging for sparse salt-and-pepper noise?
\begin{enumerate}[label=\textbf{\Alph*.}]
\item \(I_t\) is a linear filter with a fixed kernel of average weights.
\item \(I_t\) increases the magnitude of all high-frequency components.
\item \(I_t\) is less affected by extreme isolated values and can preserve edges better.
\item \(I_t\) guarantees removal of every type of noise without information loss.
\end{enumerate}
\end{minipage}
\par\vspace{12pt}
\begin{minipage}{\linewidth}
\textbf{Q12 — Multiple choice}\par\smallskip
A Gabor filter is best described as what?
\begin{enumerate}[label=\textbf{\Alph*.}]
\item A sinusoidal pattern modulated by a Gaussian envelope,\allowbreak{} selective for local frequency and orientation.
\item An unwindowed global sinusoid that has no spatial localization.
\item A median of oriented derivative responses,\allowbreak{} making it a nonlinear rank filter.
\item A Gaussian smoothing filter with no sinusoidal carrier and no orientation selectivity.
\end{enumerate}
\end{minipage}
\par\vspace{12pt}
\begin{minipage}{\linewidth}
\textbf{Q13 — Multiple choice}\par\smallskip
You multiply two same-length unpadded DFTs and take the inverse DFT. What operation does this normally implement?
\begin{enumerate}[label=\textbf{\Alph*.}]
\item Pointwise multiplication of the original spatial samples.
\item Circular convolution,\allowbreak{} with wrap-around at the signal boundaries.
\item Cross-correlation with no complex conjugation or reversal.
\item Full linear convolution of length \(N+M-1\) regardless of padding.
\end{enumerate}
\end{minipage}
\par\vspace{12pt}
\begin{minipage}{\linewidth}
\textbf{Q14 — Multiple choice}\par\smallskip
A kernel measures the x derivative of an image. Which edge orientation usually gives a large response?
\begin{enumerate}[label=\textbf{\Alph*.}]
\item Only a diagonal edge with exactly 45-degree slope.
\item A horizontal edge: intensity changes only across y.
\item A vertical edge: intensity changes across x.
\item A perfectly constant patch.
\end{enumerate}
\end{minipage}
\par\vspace{12pt}
\begin{minipage}{\linewidth}
\textbf{Q15 — Multiple choice}\par\smallskip
Compared with a small Gaussian \(\sigma\),\allowbreak{} what is a likely effect of increasing \(\sigma\) in Canny?
\begin{enumerate}[label=\textbf{\Alph*.}]
\item Suppress finer detail and noise,\allowbreak{} while potentially merging edges and reducing localization precision.
\item Remove the need for non-maximum suppression.
\item Make the operation exactly scale invariant.
\item Guarantee detection of every weak edge with perfect localization.
\end{enumerate}
\end{minipage}
\par\vspace{12pt}
\begin{minipage}{\linewidth}
\textbf{Q16 — Multiple choice}\par\smallskip
For a structure tensor at a pixel,\allowbreak{} which implementation follows the usual order of operations?
\begin{enumerate}[label=\textbf{\Alph*.}]
\item Compute \(I_x^2,\;I_x I_y\text{ and }I_y^2\) first,\allowbreak{} then aggregate or smooth those products over the neighbourhood.
\item Use \(I_{xx}\) and \(I_{yy}\) as the diagonal entries and \(I_{xy}\) off diagonal,\allowbreak{} because the structure tensor is the image Hessian.
\item Smooth \(I_x\) and \(I_y\) to one mean each,\allowbreak{} then use only the outer product of those two means.
\item Aggregate only \(I_x\) and \(I_y\) and use their signed means as the diagonal entries,\allowbreak{} setting the off-diagonal entries to zero.
\end{enumerate}
\end{minipage}
\par\vspace{12pt}
\begin{minipage}{\linewidth}
\textbf{Q17 — Multiple choice}\par\smallskip
Shi–Tomasi's corner score for structure-tensor eigenvalues \(\lambda_1\text{ and }\lambda_2\) is which expression?
\begin{enumerate}[label=\textbf{\Alph*.}]
\item \(\max(\lambda_1,\lambda_2)\)
\item \(\lambda_1+\lambda_2\),\allowbreak{} regardless of their relative sizes
\item \(\min(\lambda_1,\lambda_2)\)
\item \(\lambda_1-\lambda_2\)
\end{enumerate}
\end{minipage}
\par\vspace{12pt}
\begin{minipage}{\linewidth}
\textbf{Q18 — Multiple choice}\par\smallskip
A Lucas–Kanade window contains 100 nonzero image-gradient vectors,\allowbreak{} all parallel. Is its two-component flow uniquely observable from those equations?
\begin{enumerate}[label=\textbf{\Alph*.}]
\item No; the design matrix has rank 1 and tangential motion remains ambiguous.
\item Yes,\allowbreak{} because there are more equations than unknowns.
\item No,\allowbreak{} because least squares works only for exactly two pixels.
\item Yes,\allowbreak{} because every gradient is nonzero.
\end{enumerate}
\end{minipage}
\par\vspace{12pt}
\begin{minipage}{\linewidth}
\textbf{Q19 — Multiple choice}\par\smallskip
Move a flow estimate from a half-resolution pyramid level to full resolution. After spatially resizing the flow arrays,\allowbreak{} what else is required?
\begin{enumerate}[label=\textbf{\Alph*.}]
\item Keep the vector magnitudes unchanged because the arrays were resized.
\item Divide both displacement components by 2.
\item Multiply displacement components by 2,\allowbreak{} then warp and refine at the finer level.
\item Set every displacement to zero before the finer solve.
\end{enumerate}
\end{minipage}
\par\vspace{12pt}
\begin{minipage}{\linewidth}
\textbf{Q20 — Multiple choice}\par\smallskip
A fixed image template is tracked with SSD. The target's shape and position remain the same but all intensities gain a constant positive offset. Which claim is correct?
\begin{enumerate}[label=\textbf{\Alph*.}]
\item SSD can increase even for the true match; a mean-normalized comparison can handle an additive offset better.
\item Unnormalized correlation always removes additive and multiplicative brightness changes.
\item Only the target's geometric rotation can change SSD.
\item SSD is exactly invariant to every additive intensity change.
\end{enumerate}
\end{minipage}
\par\vspace{12pt}
\begin{minipage}{\linewidth}
\textbf{Q21 — Multiple choice}\par\smallskip
Which statement about mean-shift clustering is correct?
\begin{enumerate}[label=\textbf{\Alph*.}]
\item \(I_t\) always requires the desired number of clusters K as its central input.
\item \(I_t\) minimizes the same fixed-K objective as k-means using identical updates.
\item Its result cannot depend on the bandwidth.
\item \(I_t\) seeks density modes; kernel bandwidth influences how many modes and clusters are found.
\end{enumerate}
\end{minipage}
\par\vspace{12pt}
\begin{minipage}{\linewidth}
\textbf{Q22 — Multiple choice}\par\smallskip
Retrieve images of the same planar building facade using a query in which part of the facade is occluded by a tree. Assuming enough distinctive features remain visible,\allowbreak{} which local-keypoint approach is appropriate?
\begin{enumerate}[label=\textbf{\Alph*.}]
\item Reject every database image unless every query descriptor has a correct facade match,\allowbreak{} including descriptors on the tree.
\item Accept every nearest descriptor match as geometrically correct,\allowbreak{} because SIFT's scale robustness guarantees the same scene point.
\item Rank only by each database image's total keypoint count; descriptor appearance and geometric consistency are unnecessary.
\item Match visible local descriptors and verify a consistent homography; the occluded parts need not match.
\end{enumerate}
\end{minipage}
\par\vspace{12pt}
\begin{minipage}{\linewidth}
\textbf{Q23 — Multiple choice}\par\smallskip
What is the generic minimum data requirement for an unconstrained 2D affine transformation?
\begin{enumerate}[label=\textbf{\Alph*.}]
\item Two point pairs because an affine transformation is only rotation plus translation.
\item Four collinear point pairs always uniquely determine it.
\item One point pair because affine transforms have two parameters.
\item Three point pairs whose source points are noncollinear.
\end{enumerate}
\end{minipage}
\par\vspace{12pt}
\begin{minipage}{\linewidth}
\textbf{Q24 — Multiple choice}\par\smallskip
Which statement accurately describes backward image warping?
\begin{enumerate}[label=\textbf{\Alph*.}]
\item \(I_t\) guarantees valid source data for every destination pixel,\allowbreak{} even outside source coverage.
\item \(I_t\) forward-maps each source pixel without interpolation.
\item \(I_t\) always avoids aliasing during downsampling,\allowbreak{} even without a low-pass filter.
\item For each destination pixel,\allowbreak{} inverse-map its coordinate and sample/\allowbreak{}interpolate the source; out-of-domain samples still need a boundary rule.
\end{enumerate}
\end{minipage}
\par\vspace{12pt}
\begin{minipage}{\linewidth}
\textbf{Q25 — Multiple choice}\par\smallskip
What does the SVM kernel trick allow?
\begin{enumerate}[label=\textbf{\Alph*.}]
\item Evaluate inner products in a transformed feature space without explicitly constructing all transformed features.
\item Make every nonlinear decision boundary linear in the original coordinates.
\item Change the distance metric only after fitting a linear SVM,\allowbreak{} without changing any training computation.
\item Automatically eliminate every mislabeled training example,\allowbreak{} making soft-margin penalties unnecessary.
\end{enumerate}
\end{minipage}
\par\vspace{12pt}
\begin{minipage}{\linewidth}
\textbf{Q26 — Multiple choice}\par\smallskip
At a fixed detection IoU threshold,\allowbreak{} two same-class detections both overlap the same ground-truth object sufficiently. In score-ordered,\allowbreak{} one-to-one matching,\allowbreak{} how are they counted?
\begin{enumerate}[label=\textbf{\Alph*.}]
\item Both are TPs but each receives half a recall increment,\allowbreak{} making this identical to a single detection.
\item The first can be a TP; the second is a duplicate FP unless it can validly match another unmatched object.
\item The first is a TP and the second is ignored automatically in AP,\allowbreak{} so duplicate boxes cannot reduce precision.
\item Both are TPs because their IoUs are high,\allowbreak{} even though they claim the same object.
\end{enumerate}
\end{minipage}
\par\vspace{12pt}
\begin{minipage}{\linewidth}
\textbf{Q27 — Multiple choice}\par\smallskip
Three same-class boxes have descending scores \(A>B>C\). \(\operatorname{IoU}(A,B)>\text{threshold}\),\allowbreak{} \(\operatorname{IoU}(B,C)>\text{threshold}\),\allowbreak{} but \(\operatorname{IoU}(A,C)<\text{threshold}\). What does greedy hard NMS keep?
\begin{enumerate}[label=\textbf{\Alph*.}]
\item A only,\allowbreak{} because overlap is transitively inherited through B.
\item A and C.
\item B and C.
\item All three,\allowbreak{} because A and C do not overlap strongly.
\end{enumerate}
\end{minipage}
\par\vspace{12pt}
\begin{minipage}{\linewidth}
\textbf{Q28 — Multiple choice}\par\smallskip
Why is identical constant initialization of all interchangeable hidden units problematic?
\begin{enumerate}[label=\textbf{\Alph*.}]
\item Random initialization guarantees no overfitting.
\item Every nonzero constant necessarily causes exploding gradients on the first step.
\item Biases must always be random for any network to learn at all.
\item Symmetric units can receive identical activations and gradients,\allowbreak{} so they fail to learn distinct features.
\end{enumerate}
\end{minipage}
\par\vspace{12pt}
\begin{minipage}{\linewidth}
\textbf{Q29 — Multiple choice}\par\smallskip
Which statement about Adam and learning rate is correct?
\begin{enumerate}[label=\textbf{\Alph*.}]
\item Adam eliminates base learning-rate tuning because its moment estimates guarantee stable convergence for any chosen step size.
\item Adam uses gradient-moment estimates to adapt parameter updates,\allowbreak{} but still has a base learning-rate hyperparameter.
\item Adam uses only a running mean of gradients,\allowbreak{} so it is exactly SGD with momentum and has no second-moment adaptation.
\item Increasing Adam's base learning rate cannot destabilize training because all adaptive updates are independent of that rate.
\end{enumerate}
\end{minipage}
\par\vspace{12pt}
\begin{minipage}{\linewidth}
\textbf{Q30 — Multiple choice}\par\smallskip
A camera only rotates about its optical centre while viewing a static scene. Why is this unsuitable as the sole initialization for ordinary depth triangulation?
\begin{enumerate}[label=\textbf{\Alph*.}]
\item The pure-rotation homography is invertible,\allowbreak{} so its inverse already provides each point's unique metric depth.
\item There is no translation baseline and no depth-dependent parallax.
\item Higher image resolution creates a metric translation baseline even when the optical centre stays fixed.
\item A pure rotation makes all otherwise visible feature correspondences invalid.
\end{enumerate}
\end{minipage}
\par\vspace{12pt}
\clearpage\section*{Part B — Numerical values}
\begin{minipage}{\linewidth}
\textbf{Q31 — Numerical}\par\smallskip
A world-to-camera extrinsic transform is \(X_c=R X_w+t\),\allowbreak{} with
\[R=\begin{bmatrix}0&0&1\\0&1&0\\-1&0&0\end{bmatrix}\]
 and \(t=\begin{bmatrix}-3&2&1\end{bmatrix}^{\mathsf T}\) metres. Report the Euclidean distance of the camera centre from the world origin. Use \(C=-R^{\mathsf T}t\).
\par\medskip\textbf{Answer:} \rule{55mm}{0.4pt}
\end{minipage}
\par\vspace{12pt}
\begin{minipage}{\linewidth}
\textbf{Q32 — Numerical}\par\smallskip
An RGB image has per-channel means \(\begin{bmatrix}80&40&20\end{bmatrix}\). A pixel has RGB value \(\begin{bmatrix}100&30&10\end{bmatrix}\). Grey-World correction uses \(\text{target}=\frac{80+40+20}{3}\),\allowbreak{} then
\[\text{corrected}=\frac{\text{pixel}}{\text{channel means}}\cdot\text{target}.\]
 Ignore clipping and quantization. Report the corrected green value.
\par\medskip\textbf{Answer:} \rule{55mm}{0.4pt}
\end{minipage}
\par\vspace{12pt}
\begin{minipage}{\linewidth}
\textbf{Q33 — Numerical}\par\smallskip
Weighted photometric stereo minimizes \(\sum_j\bigl(w_j(L_j g-i_j)\bigr)^2\). The last light vector's magnitude also encodes its strength. Report \(g_2\).


\begin{Code}
import numpy as np
L = np.array([[1,0,0],[0,1,0],[0,0,1],[1,1,1]], dtype=float)
i = np.array([.2,.3,.6,1.2])
w = np.array([1,1,1,2], dtype=float)
\end{Code}


Provided routine: \texttt{np.linalg.lstsq(L*\allowbreak{}w[:,\allowbreak{}None],\allowbreak{} i*\allowbreak{}w,\allowbreak{} rcond=\allowbreak{}None)[0]} minimizes the stated weighted residuals.
\par\medskip\textbf{Answer:} \rule{55mm}{0.4pt}
\end{minipage}
\par\vspace{12pt}
\begin{minipage}{\linewidth}
\textbf{Q34 — Numerical}\par\smallskip
Use the simplified Phong model \(I=I_a+k_d\max(n\cdot l,0)+k_s\max(v\cdot r,0)^s\). Let \(I_a=0.1\),\allowbreak{} \(k_d=0.6\),\allowbreak{} \(n\cdot l=0.5\),\allowbreak{} \(k_s=0.8\),\allowbreak{} \(v\cdot r=0.5\),\allowbreak{} and \(s=4\). Report I.
\par\medskip\textbf{Answer:} \rule{55mm}{0.4pt}
\end{minipage}
\par\vspace{12pt}
\begin{minipage}{\linewidth}
\textbf{Q35 — Numerical}\par\smallskip
Linearly stretch intensities from input range \(\begin{bmatrix}50&150\end{bmatrix}\) to \(\begin{bmatrix}0&255\end{bmatrix}\). Report the output for input intensity 70. No clipping is needed for this value.
\par\medskip\textbf{Answer:} \rule{55mm}{0.4pt}
\end{minipage}
\par\vspace{12pt}
\begin{minipage}{\linewidth}
\textbf{Q36 — Numerical}\par\smallskip
For the binary image below,\allowbreak{} compute the number of foreground components under 4-connectivity and under 8-connectivity. Report \(\text{count}_4-\text{count}_8\) as one exact integer. Background is not counted.


\[
\begin{bmatrix}1&0&1&0&0\\0&1&0&0&1\\0&0&1&0&1\\0&0&0&1&0\\1&0&0&0&1\end{bmatrix}
\]

\par\medskip\textbf{Answer:} \rule{55mm}{0.4pt}
\end{minipage}
\par\vspace{12pt}
\begin{minipage}{\linewidth}
\textbf{Q37 — Numerical}\par\smallskip
Report the centre output of a 3 by 3 median filter on this patch. Use all nine entries and report an exact integer.


\[
\begin{bmatrix}0&7&8\\8&255&9\\10&11&0\end{bmatrix}
\]

\par\medskip\textbf{Answer:} \rule{55mm}{0.4pt}
\end{minipage}
\par\vspace{12pt}
\begin{minipage}{\linewidth}
\textbf{Q38 — Numerical}\par\smallskip
Construct a magnitude-weighted UNSIGNED HoG histogram from \(G_x=\begin{bmatrix}3&-3&0&2\end{bmatrix}\) and \(G_y=\begin{bmatrix}4&-4&5&0\end{bmatrix}\). Use angles \(\theta=\operatorname{degrees}(\operatorname{atan2}(G_y,G_x))\bmod180\). Use nine HARD bins [0,\allowbreak{}20),\allowbreak{}[20,\allowbreak{}40),\allowbreak{}...,\allowbreak{}[160,\allowbreak{}180),\allowbreak{} numbered 0..8; no interpolation and no histogram normalization. Report the value of bin 2,\allowbreak{} which covers [40,\allowbreak{}60) degrees.

Provided functions: \texttt{np.arctan2(y,\allowbreak{}x)} gives quadrant-aware angles in radians; \texttt{np.degrees(angle)} converts radians to degrees. Magnitude is \texttt{np.sqrt(Gx*\allowbreak{}Gx+Gy*\allowbreak{}Gy)}.
\par\medskip\textbf{Answer:} \rule{55mm}{0.4pt}
\end{minipage}
\par\vspace{12pt}
\begin{minipage}{\linewidth}
\textbf{Q39 — Numerical}\par\smallskip
A complex Fourier coefficient is \(F=-3-4\mathrm{i}\). Report its principal phase in degrees in the range (-180,\allowbreak{}180]. Provided function: \texttt{np.degrees(np.arctan2(F.imag,\allowbreak{} F.real))} uses quadrant-aware \(\operatorname{atan2}(y,x)\).
\par\medskip\textbf{Answer:} \rule{55mm}{0.4pt}
\end{minipage}
\par\vspace{12pt}
\begin{minipage}{\linewidth}
\textbf{Q40 — Numerical}\par\smallskip
In Lucas–Kanade,\allowbreak{} solve \(A\begin{bmatrix}u&v\end{bmatrix}^{\mathsf T}=-I_t\) by least squares and report v.


\begin{Code}
import numpy as np
A = np.array([[2,1],[1,-1],[0,2],[3,0]], dtype=float)
It = np.array([-5,-1,-3,-6], dtype=float)
\end{Code}


Provided routine: \texttt{np.linalg.lstsq(A,\allowbreak{} -It,\allowbreak{} rcond=\allowbreak{}None)[0]} returns [u,\allowbreak{}v].
\par\medskip\textbf{Answer:} \rule{55mm}{0.4pt}
\end{minipage}
\par\vspace{12pt}
\begin{minipage}{\linewidth}
\textbf{Q41 — Numerical}\par\smallskip
Report the zero-based index returned by this SSD matching code. NumPy's argmin returns the FIRST minimum in a tie.


\begin{Code}
import numpy as np
T = np.array([[1,2],[3,4]])
P = np.array([[[2,3],[4,5]],
              [[1,1],[3,5]],
              [[2,2],[2,4]]])
costs = ((P-T)**2).sum(axis=(1,2))
answer = int(np.argmin(costs))
\end{Code}

\par\medskip\textbf{Answer:} \rule{55mm}{0.4pt}
\end{minipage}
\par\vspace{12pt}
\begin{minipage}{\linewidth}
\textbf{Q42 — Numerical}\par\smallskip
Fit a 2D affine transform by least squares. Source rows are
\[\begin{bmatrix}0&0\\1&0\\0&1\\1&1\end{bmatrix};\]
 target rows are
\[\begin{bmatrix}2&-1\\4&-1\\3&2\\5&3\end{bmatrix}.\]
 Build A with rows \([x,y,1]\) and solve \(AP=\text{targets}\). Report the predicted target y coordinate for source (2,\allowbreak{}3).

Provided routine: \texttt{P=\allowbreak{}np.linalg.lstsq(A,\allowbreak{} targets,\allowbreak{} rcond=\allowbreak{}None)[0]}; prediction is \texttt{np.array([2,\allowbreak{}3,\allowbreak{}1])@P}.
\par\medskip\textbf{Answer:} \rule{55mm}{0.4pt}
\end{minipage}
\par\vspace{12pt}
\begin{minipage}{\linewidth}
\textbf{Q43 — Numerical}\par\smallskip
Rows contain [nearest\_\allowbreak{}distance,\allowbreak{} second\_\allowbreak{}nearest\_\allowbreak{}distance] from descriptor matching. Keep a match iff \(\frac{d_1}{d_2}<0.75\) (strict inequality). Report the count kept.


\[
\begin{bmatrix}0.3&0.5\\0.4&0.5\\0.6&1.0\\0.2&0.21\\0.1&0.4\end{bmatrix}
\]

\par\medskip\textbf{Answer:} \rule{55mm}{0.4pt}
\end{minipage}
\par\vspace{12pt}
\begin{minipage}{\linewidth}
\textbf{Q44 — Numerical}\par\smallskip
Run ONE assignment/\allowbreak{}update iteration of 1D k-means with initial centres \(\begin{bmatrix}0&10\end{bmatrix}\) and data \(\begin{bmatrix}0&2&3&9&11\end{bmatrix}\). Use squared Euclidean distance. Report the updated smaller centre.
\par\medskip\textbf{Answer:} \rule{55mm}{0.4pt}
\end{minipage}
\par\vspace{12pt}
\begin{minipage}{\linewidth}
\textbf{Q45 — Numerical}\par\smallskip
Rows are actual classes; columns are predicted classes. Report MACRO recall as a fraction for this matrix.


\[
\begin{bmatrix}30&10&0\\5&10&5\\0&2&8\end{bmatrix}
\]

\par\medskip\textbf{Answer:} \rule{55mm}{0.4pt}
\end{minipage}
\par\vspace{12pt}
\begin{minipage}{\linewidth}
\textbf{Q46 — Numerical}\par\smallskip
Greedy hard NMS processes one class in score order A,\allowbreak{}B,\allowbreak{}C,\allowbreak{}D (all scores distinct). Exact pairwise IoUs are below; fractions mean ordinary division. Suppress a lower-scoring box only when its IoU with a KEPT box is strictly greater than 0.5. Report the number kept.


\[
\begin{array}{c|rrrr}&A&B&C&D\\\hline A&1&\frac{3}{5}&\frac{1}{3}&\frac{1}{7}\\[6pt]B&\frac{3}{5}&1&\frac{3}{5}&\frac{1}{3}\\[6pt]C&\frac{1}{3}&\frac{3}{5}&1&\frac{3}{5}\\[6pt]D&\frac{1}{7}&\frac{1}{3}&\frac{3}{5}&1\end{array}
\]

\par\medskip\textbf{Answer:} \rule{55mm}{0.4pt}
\end{minipage}
\par\vspace{12pt}
\begin{minipage}{\linewidth}
\textbf{Q47 — Numerical}\par\smallskip
A CNN layer uses a valid stride-1 CROSS-CORRELATION,\allowbreak{} no bias,\allowbreak{}
\[I=\begin{bmatrix}1&3&5\\2&-1&4\end{bmatrix}\]
 and
\[W=\begin{bmatrix}1&0\\-1&2\end{bmatrix}.\]
 Its two output positions receive upstream loss gradients \(\begin{bmatrix}2&-1\end{bmatrix}\). Report \(\frac{\partial\mathrm{Loss}}{\partial W_{1,1}}\),\allowbreak{} with zero-based indices. The derivative sums \(\text{upstream gradient}\cdot\text{the corresponding input value}\) over BOTH output positions.
\par\medskip\textbf{Answer:} \rule{55mm}{0.4pt}
\end{minipage}
\par\vspace{12pt}
\begin{minipage}{\linewidth}
\textbf{Q48 — Numerical}\par\smallskip
A final feature map has shape \(9\times9\times8\). Global average pooling replaces each channel by its spatial mean. A fully connected 5-class head then uses one bias per class. Report the head's trainable parameter count; exclude the backbone.
\par\medskip\textbf{Answer:} \rule{55mm}{0.4pt}
\end{minipage}
\par\vspace{12pt}
\begin{minipage}{\linewidth}
\textbf{Q49 — Numerical}\par\smallskip
Use ORIGINAL YOLO v1,\allowbreak{} whose output shape is \(S\times S\times(5B+C)\): B boxes with 5 values each,\allowbreak{} and ONE shared C-class probability vector per cell. Let \(S=7,\;B=2\text{ and }C=4\). Report the total number of scalar outputs.
\par\medskip\textbf{Answer:} \rule{55mm}{0.4pt}
\end{minipage}
\par\vspace{12pt}
\begin{minipage}{\linewidth}
\textbf{Q50 — Numerical}\par\smallskip
An ideal rectified stereo pair has focal length \(f=1200\) pixels,\allowbreak{} baseline \(B=0.12\) metres,\allowbreak{} and matching horizontal coordinates \(u_L=414,\;u_R=390\). Both principal points are equal. Use disparity \(d=u_L-u_R\) and depth \(Z=\frac{fB}{d}\). Report Z in metres.
\par\medskip\textbf{Answer:} \rule{55mm}{0.4pt}
\end{minipage}
\par\vspace{12pt}

\end{document}
